% XeLaTeX can use any Mac OS X font. See the setromanfont command below.
% Input to XeLaTeX is full Unicode, so Unicode characters can be typed directly into the source.

% The next lines tell TeXShop to typeset with xelatex, and to open and save the source with Unicode encoding.

%!TEX TS-program = xelatex
%!TEX encoding = UTF-8 Unicode

\documentclass[11pt]{article}


%%%%%%%%%%%%%%%
\def\Type{ApplicationPreDJ}
\def\TypeNum{Pre-class for Application Day: Gradient Descent Method}
\def\TypeTitle{}
\def\Semester{Fall 2025}
\def\Course{Math 1190/1210}
\def\Targets{{\bf\color{Goldenrod}AD1}}

\usepackage{preambleDJ}




\begin{document}
\pagestyle{otherpages}
\thispagestyle{firstpage}
\vspace*{.65in}

%If Ada Lovelace\footnote{\url{https://en.wikipedia.org/wiki/Ada_Lovelace}} is the ``mother of programming" then 
\mpageT{.35}{.65}{
%\graphicsR{adaLovelace2}{2}
\graphicsL{feiFeiLiWiredCropped}{2}
}{
{\bf Introduction}\\
Fei-Fei Li is  known as the ``godmother of A.I" for her  groundbreaking contributions in computer vision and deep learning and as the creator of ImageNet.\footnotemark 
~\\\\
In this application, we will be exploring how the Gradient Descent method can be used for a supervised machine learning model that we have already seen in the Keeling Curve application. Here's an example.\\
\example Given $\ds z=f(x,y)=4-yx^2-y^2$, if $(x,y)=(1,2)$, $$z=f(1,2)=4-2\cdot 1^2-2^2=-2$$ This allows us to express a point in $3$D as $(x,y,z)=(1,2,-2)$.\\
{\bf Graphing Points and Surfaces} \\
\mpageT{.5}{.5}[\hs{.05}]{
Click on link to go to\\
 \url{https://www.desmos.com/3d/1bwyuvhfsp}  \\
Click on the wrench icon and check ``Reverse contrast", ``Translucent surfaces" and move\\
}{
\graphicsR{contrastTranslucentPerspective.pdf}{2.5}
}
\vs{-.2}\\
the bounding box perspective slider to the far right for better viewing.
}
\footnotetext{
See ``Godmother of A.I." Fei-Fei Li on technology development: ``The power lies within people" at \url{https://www.cbsnews.com/news/godmother-of-a-i-on-technology-development-fei-fei-li/}\\
 or Britannica at \url{https://www.britannica.com/biography/Fei-Fei-Li}, The Guardian \href{https://www.theguardian.com/technology/2023/nov/05/ai-pioneer-fei-fei-li-im-more-concerned-about- the-risks-that-are-here-and-now}{\underline{here}} or \url{https://en.wikipedia.org/wiki/Fei-Fei_Li} }

\vspace{0.5cm}

\mpageT{.7}{.4}{
The first things we see are the axes.
\itemC{
\item $x$ axis (positive direction is left)
\item $y$-axis (positive direction is towards the viewer)
\item $z$-axis (positive direction is up)
}
We also see the graph of our point $(1.2,-2)$ from Example 1.  The arrows show that we move $1$ unit to the left from the origin (along the $x$-axis), then $2$ units towards the viewer (parallel to the $y$-axis) and then $2$ units down (parallel to the $z$-axis and down because $z$-value is negative).
}{
\graphicsR{pointInDesmos}{2.5}
}~\\

\newpage
An infinite collection of these points form a surface. Here are some examples of surfaces.
\def\sz{1.}
\tasksA{4}{
\task A sphere\vs{-.3}\\
\graphicsL{sphere.png}{\sz}
\task A flat plane\vs{-.3}\\
\graphicsL{plane}{\sz}
\task A saddle\vs{-.3}\\
\graphicsL{saddle.png}{\sz}
\task A {\em paraboloid} (i.e. a bowl.)\vs{-.25}\\
\graphicsL{parabaloid.png}{\sz}
%\task* I have no idea, because it's a big blob of points. ({\em If you are considering this choice, you can get a better idea of the surface by clicking on the grey rectangle to see the underlying surface.  Clicking on white circle removes all the points.})
}

{\bf Derivatives of Functions of 2 Variables} \\
\mpageT{.6}{.4}{
Go to \url{https://www.desmos.com/3d/0c2zwfghdx} \hfill \\
Click on the wrench icon and check ``Reverse contrast", ``Translucent surfaces" and move the bounding box perspective slider to the far right for better viewing.
}{
\graphicsR{contrastTranslucentPerspective.pdf}{2.5}
}~\\\\
 Since there are two inputs in a function of the form $z=f(x,y)$, there are {\em two} derivatives.  We will interpret each as we have interpreted derivatives previously: ``the derivative is the slope of a tangent line at the given point". Desmos will help us visualize and interpret each of these derivatives.\\\\
 {\bf \em Interpreting Derivatives Graphically}\\
 In Desmos, you will see the graph of the paraboloid $z=f(x,y)=4-yx^2-y^2$ with the point $P=(x,y,f(x,y))=(1,2,-2)$.  
\enumR{
\item To find the derivative of $f(x,y)$ with respect to $x$ at $(1,2)$.  
\itemI{
\item Click on the red rectangle to toggle on/off the graphics related to the derivative of $f$ with respect to  $x$.
\item Since $(x,y)=(1,2)$ the vertical plane $y=2$ appears and intersects the  surface to create an upside down parabola parallel to the $x$-axis.  Rotate the object to better see this parabola via click/dragging the graph.
\item You will also see a tangent line parallel to the $x$-axis through the point $P=(1,2,-2)$. {\em The slope of this tangent line is the derivative of $f$ with respect to $x$ at $(1,2)$.}
\item {\em Notation}:  ``The derivative of $f$ with respect to $x$ at $(1,2)$" can be denoted two different (but equivalent) ways.
\itemI{
\item $\ds \frac{d}{dx} f(x,y)\Big|_{(1,2)}$
\item $\ds f_x(1,2)$
}
\item Use Desmos to find $\ds  f_x(1,2)$. {\em Hint}: Look for $\ds f_x(x_{point},y_{point})$.\\
% $\ds  f_x(1,2)=$\unline{.5}{\blue{$-4$}}
}
\item To find the derivative of $f(x,y)$ with respect to $y$ at $(1,2)$.  
\itemI{
\item Click on the blue rectangle to toggle on/off the graphics related to the derivative of $f$ with respect to  $y$.
\item Since $(x,y)=(1,2)$ the vertical plane $x=1$ appears and intersects the  surface to create an upside down parabola parallel to the $y$-axis.  Rotate the object to better see this parabola via click/dragging the graph.
\item You will also see a tangent line parallel to the $y$-axis through the point $P=(0,-1/2,3/4)$. {\em The slope of this tangent line is the derivative of $f$ with respect to $y$ at $(1,2)$}
\item {\em Notation}:  ``The derivative of $f$ with respect to $y$ at $(1,2)$" can be denoted two different (but equivalent) ways.
\itemI{
\item $\ds \frac{d}{dy} f(x,y)\Big|_{(1,2)}$
\item $\ds f_y(1,2)$
}
\item Use Desmos to find $\ds  f_y(1,2)$. {\em Hint}: Look for $\ds f_y(x_{point},y_{point})$.\\
% $\ds  f_y(1,2)=$\unline{.5}{\blue{$-5$}}
}
}
The locations where $f_x(x,y)=0$ and $f_y(x,y)=0$ at the same time are called {\em critical points}.  They indicate the possibility of  a special property at that location. In our example, we're at a {\em saddle point}. In other instances, we might be at a {\em local maximum (i.e. a peak)} or a {\em local minimum (i.e. a valley)}.
~\\\\
{\bf \em Interpreting Derivatives Symbolically}\\
Our visual interpretation of derivatives suggests an algorithm to find function equations of derivatives.     When we find symbolic derivatives with respect to $x$, we're going to do the same-- assume $y$ is some constant-- and then use our derivative shortcuts to find $f_x(x,y)$.\\\\
\example Given $f(x,y)=4-yx^2-y^2$, find $\ds f_x(1,2)$ symbolically.  
\itemI{
\item When we found $f_x(1,2)$ graphically in Desmos, we set $y=2$ to create a plane that sliced through our surface. This allowed us to visualize a curve and tangent line.
\item To find $\ds f_x(1,2)$ symbolically, we again assume $y$ is some constant (it might be $y=2$ or any other constant).
\item Since $f(x,y)=4-yx^2-y^2$,
\al{
f_x(x,y)&=\frac{d}{dx}4-\frac{d}{dx}yx^2-\frac{d}{dx}y^2\\
f_x(x,y)&=\frac{d}{dx}4-y\frac{d}{dx}x^2-\frac{d}{dx}y^2 \qquad\text{\em \footnotesize  (Since $y$ is a constant, we can bring it out front of $\ds \frac{d}{dx}x^2$ }\\
f_x(x,y)&=0-y\cdot 2x-0\qquad\text{\em \footnotesize  (Since $y$ is a constant, $(y-1)^2$ is constant and thus $\ds \frac{d}{dx}(y-1)^2=0$)}\\
f_x(x,y)&=-2xy
\intertext{We now have a symbolic equation for $f_x$ that we can evaluate at any point $(x,y)$. We use the point we've been dealing with: $(x,y)=(1,2)$.}
f_x(1,2)&=-2\cdot 1 \cdot 2\\
&=-4
}
}


\example Given $f(x,y)=4-yx^2-y^2$, find $\ds f_y(1,2)$ symbolically. 
\itemI{
\item When we found $f_x(1,2)$ graphically in Desmos, we set $x=1$ to create a plane that sliced through our surface. This allowed us to visualize a second curve and tangent line.
\item To find $\ds f_y(1,2)$ symbolically, we again assume $x$ is some constant (it might be $x=1$ or any other constant).
\item Since $f(x,y)=4-yx^2-y^2$,
\al{
f_y(x,y)&=\frac{d}{dy}4-\frac{d}{dy}yx^2-\frac{d}{dy}y^2\\
f_y(x,y)&=\frac{d}{dy}4-x^2\frac{d}{dy}y-\frac{d}{dy}y^2\qquad\text{\em \footnotesize  (Since $x$ is a constant, $x^2$ is constant and we can move it in front of $\ds \frac{d}{dy} y$. }\\
&=0-x^2\cdot 1-2y \\
&=-x^2-2y\\
f_y(1,2)&=-(1)^2-2(2)\\
&=-5
}
}

{\bf \em Summary: Finding Derivatives Symbolically}
\itemI{
\item To find the derivative of $f$ with respect to $x$, $f_x(x,y)$, assume $y$ is constant and find the derivative with respect to $x$ using our usual derivative shortcuts.
\item To find the derivative of $f$ with respect to $y$, $f_y(x,y)$, assume $x$ is constant and find the derivative with respect to $y$ using our usual derivative shortcuts.
}
~\\

%~\\\\
% \q{Should we ask this question} Given $h(x,y)=e^{-(x^2-xy+y^2)}$,  find $h_x(x,y)$ and $h_y(x,y)$. Note that your answers will contain both $x$ and $y$.
%{\em Hint:} Don't forget to use the Chain Rule! 
%\enum{
%\item $h_x(x,y)=$\unline{.75}{\blue{$-(2x-y)e^{-(x^2-xy+y^2)}$}}
%\item $h_y(x,y)=$\unline{.75}{\blue{$-(-x+2y)e^{-(x^2-xy+y^2)}$}}
%}
%}


{\bf Vectors and the Gradient} \\
It is common to group our two derivatives ($f_x$ and $f_y$) together to create a {\em vector} that we call the {\em gradient}. In particular the notation we use is $\nabla f(x,y)=<f_x(x,y),f_y(x,y)>$  In our previous example we saw that at the point $(1,2)$, $$\nabla f(1,2)=<f_x(1,2),f_y(1,2)>=<-4,-5>.$$ We will explore what a general vector is and then come back to considering  special properties of a gradient vector.\\

{\bf \em Vector Properties} Our vectors will have two components and be in the form $<x,y>$ where  the $x$ component  describes the change in the $x$-direction and $y$ component describes the change in the $y$direction.  We graph vectors as an arrow. Each vector has two properties: {\em length} and {\em direction}.  \\
\mpageT{.65}{.35}{
{\em Direction of a vector}
\itemI{
\item To determine direction, we graph $<-4,-5>$ by starting at a base point $(0,0)$ in Figure 1) and  move $2$ units to the left and $3$ units down. The end point of the vector is at the point $(-4,-5)$.  See Figure 1.\\\\
Note that location is {\em not} a vector property. This means that a vector can have any base point. If two vectors have identical direction and length, then they are considered the same vector. \\\\}
%
}{
\graphicsR{vector1}{2}{\bf Figure 1}}

\mpageT{.65}{.35}{
Figure 2 shows the  vector $<3,4>$ in two different locations. In the first quadrant $<5,4>$  has a base point of $(0,0)$ and an endpoint of $(3,4)$\\\\
\example The vector $<3,4>$ in the third quadrant has a base point of $(-4,-5)$.  We add the vector components to the base point components to find the end point.  
get $$(x?,y?)=(-4+3,-5+4)=(-1,-1)$$ 
} 
{
\graphicsR{vector2}{2}{\bf Figure 2}
}

\mpageT{.65}{.35}{
{\em Length of a vector}
\itemI{
\item To determine length, we can interpret the vector as the hypotenuse of a triangle with one side length $|-4|=4$ and the other side a length of $|-5|=5$. We use Pythagorean's theorem (Figure 3) to get 
$$\text{Length}= |<-4,-5>|= \sqrt{(-4)^2+(-5)^2}=\sqrt{41}\approx6.40$$
Note that the notation $ |<-4,-5>|$ means {\em ``find the length of $<-4,-5>$"}\\\\
}
}{
\graphicsR{triangle}{2}{\bf Figure 3}
}~\\
\example
Find  $|<3,4>|$.  Recall that Figure 2 show the graph of this vector in two locations.  $$|<3,4>|=\sqrt{3^2+4^2}=\sqrt{9+16}=\sqrt{25}=5$$.
%{\bf During Class}\\
%During class we will see how Gradient Vectors are used in a machine learning context.  We're going to see that gradient vectors act like compass and point in the direction where a surface is increasing the fastest. \\

%More specifically, in the Keeling Curve application we created lines of best fit of the form $$c_1\sim mt_1+b$$ where Desmos provided values for $m$ and $b$. We're going to examine a machine learning algorithm to estimate these parameters using Gradient Vectors.   
\end{document}  
